El aseguramiento de calidad del proyecto se ha abordado como un conjunto de prácticas técnicas y documentales orientadas a conseguir un sistema mantenible, comprensible, verificable y coherente con los requisitos definidos. Dado que se trata de una aplicación web \textit{full-stack} desarrollada de forma iterativa, la calidad no se ha concentrado únicamente en una fase final, sino que se ha revisado de forma progresiva durante el análisis, el diseño, la construcción, las pruebas, el despliegue de demostración y la elaboración de la memoria.

Como marco general, se han tomado como referencia las características de calidad recogidas por ISO/IEC 25010, especialmente mantenibilidad, seguridad, adecuación funcional, compatibilidad, usabilidad y fiabilidad \citep{iso25010}. A partir de esta referencia, el proyecto ha concretado sus controles de calidad en las siguientes prácticas.

\textbf{Principios de \textit{Clean Code}:} Durante el desarrollo se han aplicado criterios inspirados en \textit{Clean Code} \citep{martin2008cleancode}, priorizando nombres descriptivos, funciones con responsabilidades acotadas, reducción de duplicidades y organización clara del código. Esta práctica ha sido especialmente relevante en servicios de dominio, validaciones de entrada, contratos compartidos, utilidades de catálogo, pedidos, visitas, comunicaciones y operaciones administrativas. El objetivo ha sido facilitar la comprensión del sistema y reducir el coste de introducir cambios posteriores.

\textbf{Separación de responsabilidades y arquitectura mantenible:} La aplicación no se ha planteado como una arquitectura hexagonal estricta, pero sí ha seguido una separación práctica de responsabilidades inspirada en los principios de \textit{Clean Architecture} \citep{martin2017cleanarchitecture}. La interfaz se organiza mediante rutas y componentes de \textit{Next.js}; la frontera de servidor se apoya en \textit{Route Handlers}; la lógica de negocio se concentra en servicios reutilizables; y la persistencia se gestiona mediante entidades, repositorios y migraciones de \textit{TypeORM}. Esta organización reduce el acoplamiento entre interfaz, reglas de negocio y base de datos, y facilita la validación independiente de partes críticas del sistema.

\textbf{Tipado, análisis estático y construcción reproducible:} El proyecto utiliza \textit{TypeScript} como mecanismo de prevención temprana de errores en contratos, componentes y servicios \citep{typescriptdocs2026}. Además, se emplea \textit{ESLint} para detectar patrones problemáticos y mantener una estructura de código consistente \citep{eslintdocs2026}. La construcción completa con \texttt{npm run build} permite comprobar que la aplicación puede compilarse correctamente con \textit{Next.js}, incluyendo rutas, componentes de servidor, componentes de cliente, configuración de imágenes y dependencias externas.

\textbf{Testing y validación progresiva:} La validación del sistema combina comprobaciones automáticas, pruebas funcionales por bloques de negocio y revisiones manuales de flujos críticos. Aunque la versión actual no incorpora todavía una batería formal de pruebas unitarias con un \textit{runner} dedicado, sí dispone de scripts específicos para comprobar servicios, persistencia, migraciones, permisos, rutas, operaciones de pedido, comunicaciones, auditoría, \textit{rate limiting} y preparación de datos de demostración. Esta estrategia permite detectar regresiones tras cada cambio relevante y aporta evidencias repetibles sobre el comportamiento del sistema.

\textbf{Calidad de interfaz, accesibilidad y experiencia de uso:} Además de las comprobaciones funcionales, se han utilizado herramientas auxiliares como \textit{React Doctor} para revisar patrones de mantenibilidad y rendimiento en componentes React \citep{reactdoctor2026}, y \textit{WAVE} como apoyo a la revisión de accesibilidad web \citep{wavewebaim2026}. Estas comprobaciones no sustituyen una auditoría profesional completa de accesibilidad, pero ayudan a detectar problemas comunes y a reforzar la consistencia de la interfaz en escritorio y móvil.

\subsection{Verificación, validación y control de calidad}

Durante el desarrollo se han aplicado varias actividades de verificación, validación y control de calidad con el objetivo de comprobar que el producto construido se mantiene alineado con el alcance definido.

\textbf{Revisión de código:} Al tratarse de un desarrollo individual, no se han realizado revisiones cruzadas formales entre varios desarrolladores. En su lugar, se ha aplicado una revisión incremental después de cada cambio relevante, comprobando legibilidad, coherencia con la arquitectura existente, ausencia de duplicidades evidentes, mantenimiento de permisos por rol y compatibilidad con los flujos ya validados. Esta revisión ha permitido corregir problemas de organización interna, consolidar servicios reutilizables y evitar que la documentación quedase desalineada respecto al estado real del repositorio.

\textbf{Comprobaciones estáticas:} Antes de considerar estable una versión, se ejecutan comprobaciones de tipado, análisis estático y construcción:

\begin{lstlisting}[language=bash]
npm run typecheck
npm run lint
npm run build
\end{lstlisting}

Estas comprobaciones permiten detectar errores de tipos, importaciones incorrectas, problemas de configuración, fallos de compilación y usos no válidos de componentes o APIs del framework.

\textbf{Pruebas funcionales y de integración:} Los flujos principales se validan mediante scripts específicos que ejercitan servicios de negocio, persistencia, migraciones y reglas de acceso. Entre los comandos utilizados se encuentran:

\begin{lstlisting}[language=bash]
npm run test:business
npm run test:modules
npm run test:all
\end{lstlisting}

El comando \texttt{npm run test:all} agrupa la comprobación estática, la ejecución de migraciones, las pruebas funcionales por áreas del sistema y las comprobaciones auxiliares de calidad de interfaz. De este modo, se obtiene una señal amplia sobre el estado técnico del proyecto antes de continuar con nuevas iteraciones.

\textbf{Pruebas de sistema y validación manual:} Además de los scripts automáticos, se han realizado validaciones manuales sobre los recorridos de administrador, comercial y cliente profesional. Estas comprobaciones incluyen autenticación, autorización por rol, navegación protegida, gestión de usuarios, clientes, visitas, rutas, catálogo, pedidos, pagos parciales, comunicaciones, promociones, formaciones, notificaciones, auditoría y operaciones de despliegue. La validación manual resulta necesaria porque algunos aspectos de experiencia de usuario, claridad visual, comportamiento responsive y adecuación al proceso real del distribuidor no pueden evaluarse completamente mediante scripts.

\textbf{Criterios de aceptación y rechazo:} Una funcionalidad se considera aceptada cuando cumple los requisitos asociados, supera las comprobaciones estáticas, no rompe flujos ya validados, mantiene la estructura técnica del proyecto, respeta los permisos definidos por rol y queda reflejada en la documentación correspondiente cuando afecta al alcance, al diseño, a la construcción o a las pruebas. Si alguno de estos puntos no se cumple, la funcionalidad se considera pendiente de ajuste antes de formar parte de una versión estable.

\textbf{Procedimientos correctivos y preventivos:} Cuando se detecta una desviación, error o inconsistencia, se aplican acciones correctoras sobre el código, la configuración, los datos de prueba o la documentación. Además, se revisa la causa del problema para evitar su repetición en iteraciones posteriores. Ejemplos de estas acciones han sido la adaptación del despliegue a \textit{Netlify}, la reubicación del \textit{rate limiting} en \textit{Route Handlers}, la consolidación de pruebas por áreas funcionales, la corrección de rutas comerciales y la actualización de la memoria conforme evolucionaba el proyecto.

Con estas prácticas, el aseguramiento de calidad del proyecto busca garantizar que el sistema entregado sea funcional dentro del alcance definido, mantenible para futuras ampliaciones y suficientemente verificable para justificar las decisiones tomadas durante el desarrollo del \textit{Trabajo Fin de Grado}.
